\documentclass[12pt,a4paper]{article}
\usepackage[a4paper,margin=1.8cm]{geometry}
\usepackage[T2A]{fontenc}
\usepackage[utf8]{inputenc}
\usepackage[russian]{babel}
\usepackage{amsmath,amssymb,mathtools}
\usepackage{array,booktabs,longtable}
\usepackage{xcolor}
\usepackage{hyperref}
\usepackage{tikz}
\hypersetup{colorlinks=true,linkcolor=blue!55!black,urlcolor=blue!55!black}
\setlength{\parindent}{0pt}
\setlength{\parskip}{5pt}
\renewcommand{\arraystretch}{1.18}
\newcommand{\err}{\textcolor{red}{Ошибка:}}
\newcommand{\idea}{\textcolor{blue!60!black}{Ключевая идея:}}
\newcommand{\result}{\textcolor{teal!60!black}{Итог:}}
\newcommand{\tasktitle}[1]{\phantomsection\label{task:#1}\section*{Задание №#1}}

\begin{document}

\hypersetup{pageanchor=false}
\begin{titlepage}
\centering
\vspace*{0.22\textheight}
{\LARGE\bfseries Ревью домашней работы\par}
\vspace{1.6cm}
{\large Автор: Лёвин Артём Александрович\par}
\vspace{0.5cm}
{\large 3 октября 2026 г.\par}
\vfill
\begin{tikzpicture}[scale=1]
\draw[blue!45!black, line width=0.7pt] (0,0) .. controls (2,0.65) and (4,-0.65) .. (6,0);
\end{tikzpicture}
\end{titlepage}
\hypersetup{pageanchor=true}

\newpage
\section*{Сводная таблица}
\small
\begin{longtable}{>{\centering\arraybackslash}p{0.8cm} p{3.1cm} p{5.2cm} p{2.5cm} p{2.8cm}}
\toprule
№ & Тема & Суть ошибки & Ваш ответ & Правильный ответ \\
\midrule
\endfirsthead
\toprule
№ & Тема & Суть ошибки & Ваш ответ & Правильный ответ \\
\midrule
\endhead
\hyperref[task:1]{1} & Дробное неравенство с модулями & Умножение на знаменатель выполнено без определения его знака. Из-за этого допустимая область решения выбрана неверно. & $[-6;-4]$ & $(-\infty;-8)\cup(-4;+\infty)$ \\
\bottomrule
\end{longtable}
\normalsize

\newpage
\vspace*{0.32\textheight}
\begin{center}
{\LARGE\bfseries Разбор ошибок}
\end{center}
\vfill
\newpage

\tasktitle{1}

\textbf{Текст задания.}
Решите неравенство
\[
\frac{|x+4|}{|x+6|-2}\ge 1.
\]

\textbf{Ваше решение.}
В начале записано исходное неравенство, после чего выполнен переход
\[
|x+4|\ge |x+6|-2.
\]
Далее выражение с модулем разбивается на случаи, и в конце получен ответ
\[
x\in[-6;-4].
\]

\textbf{Ваш ответ.}
\[
[-6;-4].
\]

\err{} знаменатель перенесён умножением без предварительного определения его знака.

\textbf{Где именно возникает ошибка.}
Исходная запись
\[
\frac{|x+4|}{|x+6|-2}\ge 1
\]
верна. Первый достоверно неверный переход ---
\[
|x+4|\ge |x+6|-2.
\]
Так можно было бы сделать только в тех точках, где знаменатель положителен. Если знаменатель отрицателен, при умножении знак неравенства должен измениться. Если знаменатель равен нулю, исходная дробь вообще не определена.

\textbf{Почему это неверно.}
Знаменатель равен
\[
|x+6|-2.
\]
Сначала нужно определить, где он положителен, отрицателен или равен нулю.

Равенство нулю возникает при
\[
|x+6|=2,
\]
то есть при
\[
x=-8\quad\text{или}\quad x=-4.
\]
Эти два значения не входят в область допустимых значений.

Кроме того, числитель
\[
|x+4|
\]
всегда неотрицателен. Поэтому если знаменатель отрицателен, вся дробь неположительна и не может быть не меньше единицы. Следовательно, для решения исходного неравенства знаменатель обязан быть положительным.

\idea{} в дробном неравенстве нельзя просто умножать обе части на выражение с переменной. Сначала нужно установить знак этого выражения. Здесь это особенно удобно: числитель неотрицателен, поэтому отрицательный знаменатель сразу исключает возможность получить значение дроби не меньше единицы.

\textbf{Исправление от последнего правильного шага.}
Требуем положительность знаменателя:
\[
|x+6|-2>0.
\]
Отсюда
\[
|x+6|>2,
\]
поэтому
\[
x<-8\quad\text{или}\quad x>-4.
\]

Только на этих промежутках знаменатель положителен, поэтому теперь можно безопасно умножить обе части исходного неравенства на него:
\[
|x+4|\ge |x+6|-2.
\]

Дальше проверяем два допустимых промежутка.

Если
\[
x<-8,
\]
то
\[
|x+4|=-x-4,
\qquad
|x+6|=-x-6.
\]
Получаем
\[
-x-4\ge -x-6-2,
\]
то есть
\[
-x-4\ge -x-8.
\]
После сокращения переменной остаётся истинное неравенство
\[
-4\ge -8.
\]
Значит, подходит весь промежуток
\[
(-\infty;-8).
\]

Если
\[
x>-4,
\]
то одновременно
\[
|x+4|=x+4,
\qquad
|x+6|=x+6.
\]
Получаем
\[
x+4\ge x+6-2,
\]
то есть
\[
x+4\ge x+4.
\]
Это тождество, поэтому подходит весь промежуток
\[
(-4;+\infty).
\]

\textbf{Полное правильное решение.}
Рассмотрим неравенство
\[
\frac{|x+4|}{|x+6|-2}\ge1.
\]

Сначала найдём область допустимых значений:
\[
|x+6|-2\ne0.
\]
Следовательно,
\[
|x+6|\ne2,
\]
откуда
\[
x\ne-8,
\qquad
x\ne-4.
\]

Так как числитель неотрицателен, при отрицательном знаменателе дробь будет неположительной. Она не сможет быть не меньше единицы. Поэтому необходимо
\[
|x+6|-2>0.
\]
Решаем:
\[
|x+6|>2,
\]
\[
x<-8\quad\text{или}\quad x>-4.
\]

На этих промежутках знаменатель положителен, поэтому умножение не меняет знак неравенства:
\[
|x+4|\ge|x+6|-2.
\]

Для $x<-8$ получаем
\[
-x-4\ge -x-8,
\]
что верно для каждого такого $x$.

Для $x>-4$ получаем
\[
x+4\ge x+4,
\]
что также верно для каждого такого $x$.

Следовательно,
\[
x\in(-\infty;-8)\cup(-4;+\infty).
\]

\textbf{Проверка.}
Возьмём точку $x=-9$ из первого промежутка. Тогда
\[
\frac{|-9+4|}{|-9+6|-2}
=
\frac{5}{3-2}
=5,
\]
а пять не меньше единицы.

Возьмём точку $x=0$ из второго промежутка. Тогда
\[
\frac{|4|}{|6|-2}
=
\frac{4}{4}
=1.
\]

Граничные значения $x=-8$ и $x=-4$ исключаются, потому что при них знаменатель равен нулю.

Для сравнения, точка $x=-6$, которая попала в Ваш ответ, даёт
\[
\frac{|-2|}{|0|-2}
=
\frac{2}{-2}
=-1,
\]
поэтому исходное неравенство в этой точке не выполняется.

\result{}
\[
(-\infty;-8)\cup(-4;+\infty).
\]

\subsection*{Задача на закрепление}
Решите неравенство
\[
\frac{|x-1|}{|x+3|-2}\ge1.
\]

\vfill
\begin{center}
\small Лёвин Артём Александрович эксклюзивно для Анны Трапезниковой
\end{center}

\end{document}
